\section{Branch prediction (\texttt{bpu})}
\textit{Sources: \texttt{design/core/bpu/src/\{bpu,bht,btb,ras\}.sv} and its
wiring in \texttt{core\_top.sv}.}

A branch is a problem for a pipeline. The core does not know if a branch is taken
until the IE stage. Without help, it must wait. The branch prediction unit (BPU)
guesses the branch early, so the pipeline can keep fetching. If the guess is
wrong, the core fixes it later and throws away the wrong path.

\diagram{bpu}{The BPU. It is read at the ID stage to guess the branch, and
updated at the IE stage when the real result is known. Three tables work
together: the BHT guesses taken or not-taken, the BTB holds the target, the RAS
predicts returns.}{fig:bpu}{0.95}

\subsection{The three tables}
\textbf{BHT (branch history table).} This guesses the \emph{direction}: taken or
not-taken. It has 1024 entries. Each entry is a 2-bit counter (strong/weak,
taken/not-taken). It uses ``gshare'': the index is the PC mixed (XOR) with an
8-bit record of the last 8 branch results. This lets the same branch get a
different guess depending on recent history.

\textbf{BTB (branch target buffer).} This holds the \emph{target} address. It has
128 entries, direct-mapped, with a tag to confirm a hit. Each entry also says if
the branch is unconditional (a jump, always taken) and if it is compressed
(16-bit). A jump is predicted taken without asking the BHT.

\textbf{RAS (return address stack).} This predicts function returns
(\sig{jalr} back to the caller). It is an 8-deep stack. A call pushes the return
address; a return would pop it.

\subsection{When it runs}
The BPU is \emph{read} at the ID stage using \sig{pc\_id}. If the BTB hits and
the branch is predicted taken, the core redirects fetch to the target. The BPU is
\emph{updated} at the IE stage, where the real branch result is known: the BHT
counter moves toward the result, and the BTB learns the target.

\subsection{Spotting a wrong guess}
At the IE stage the core compares the real result with the guess. A guess is
wrong if the direction is wrong, or if the direction was right but the target was
wrong. On a wrong guess the core redirects to the correct PC and flushes the
wrong-path instructions in IF/ID.

\subsection{What is on and off today}
Not every BPU path is enabled. The code keeps some off on purpose:
\begin{center}\small
\begin{tabularx}{\textwidth}{l l Y}
\toprule
\textbf{path} & \textbf{status} & \textbf{why} \\
\midrule
BHT + BTB at ID & \textbf{on}  & the main predictor; does most of the work \\
RAS push        & \textbf{on}  & keeps the return stack correct; no redirect, so safe \\
RAS pop redirect& \textbf{off} & it broke Linux boot: the pop-redirect disturbed ITLB/PTW state and made the kernel's first \sig{satp} write page-fault \\
IF-stage predict& \textbf{off} & costs about 4\% on Dhrystone, but ID-stage prediction already gives most of the gain; kept off for simpler debug and Linux \\
\bottomrule
\end{tabularx}
\end{center}
The RAS still pushes and pops internally (so its state stays correct), it just
does not drive a fetch redirect.
